% Generated draft: requires mathematical and source review.
\documentclass[11pt]{article}
\usepackage[margin=25mm]{geometry}
\usepackage{fontspec}
\setmainfont{Latin Modern Roman}
\usepackage{amsmath,amssymb,mathtools,unicode-math}
\setmathfont{Latin Modern Math}
\usepackage{enumitem,needspace,xurl,hyperref}
\setlength{\parindent}{0pt}
\setlength{\parskip}{5pt}
\setlength{\emergencystretch}{3em}
\newcommand{\N}{\mathbb{N}}
\newcommand{\Z}{\mathbb{Z}}
\newcommand{\Q}{\mathbb{Q}}
\newcommand{\R}{\mathbb{R}}
\newcommand{\C}{\mathbb{C}}
\newcommand{\F}{\mathbb{F}}
\newcommand{\A}{\mathbb{A}}
\newcommand{\PP}{\mathbb{P}}
\newcommand{\E}{\mathbb{E}}
\newcommand{\eps}{\varepsilon}
\newcommand{\dd}{\,\mathrm{d}}
\begin{document}
\section*{UnsolvedMath ID: 701077}
\textbf{Source problem code:} catalogue entry
\subsection{Definitions and mathematical statement}
Let \(p\) be a prime and let \(\F_p=\Z/p\Z\). Under the standard interpretation of the Nottingham group, put
\[
N(\F_p)=\left\{f(x)=x+\sum_{n\geq 2}a_nx^n : a_n\in\F_p\right\},
\]
with group operation \((f\circ g)(x)=f(g(x))\) and identity \(x\). Give this group the topology whose neighbourhoods of the identity are
\[
U_m=\{f\in N(\F_p):f(x)\equiv x\pmod{x^{m+1}}\},\qquad m\geq 1.
\]
This is the usual pro-\(p\) topology. The supplied statement instead refers to Archive Problem 12.24 for its definition; that referenced text is absent from the supplied record, so the identification with this standard definition requires source review.

\textbf{Problem.} Find defining relations for \(N(\F_p)\) as a pro-\(p\) group. More precisely, find a topological generating set \(S\) and a set \(R\) of elements of the free pro-\(p\) group \(F_{\mathrm{pro}\text{-}p}(S)\) such that the natural continuous homomorphism onto \(N(\F_p)\) has kernel equal to the closed normal subgroup generated by \(R\). Thus the requested presentation is
\[
N(\F_p)\cong F_{\mathrm{pro}\text{-}p}(S)/\overline{\langle\!\langle R\rangle\!\rangle}.
\]
No finiteness requirement on \(R\) is stated in the source.

The source asserts that the group has two generators, for example
\[
u(x)=x+x^2,\qquad v(x)=\frac{x}{1-x}=\sum_{n\geq 1}x^n.
\]
Here the fraction denotes a formal power series. In the asserted two-generator formulation, let \(F=F_{\mathrm{pro}\text{-}p}(a,b)\) and let \(\pi:F\to N(\F_p)\) be the continuous homomorphism with \(\pi(a)=u\) and \(\pi(b)=v\). Subject to the source's assertion that these elements topologically generate the group, the task is to describe \(R\subseteq F\) with
\[
\ker\pi=\overline{\langle\!\langle R\rangle\!\rangle}_F.
\]
The supplied wording does not restrict the prime in its generator assertion; its validity and intended prime range require verification, including at \(p=2\).

\textit{Supplied qualification.} A comment dated 2009 directs the reader, for \(p>2\), to progress by M. V. Ershov in the 2005 reference below. The record does not describe the results or their effect on the remaining task.
\subsection{Short English statement}
Find a presentation of the Nottingham group over \(\F_p\) using relations in a free pro-\(p\) group, with the relation subgroup taken to be closed. The source proposes the generators \(x+x^2\) and \(x/(1-x)\), but their stated prime range needs review. Its 2009 comment cites progress for \(p>2\) without specifying the results.
\subsection{Sources}
\begin{itemize}
\item Local UnsolvedMath registry, supplied complete source record: permanent ID \(701077\), exact source code \texttt{catalogue entry}, title \emph{Kourovka Notebook Problem 14.93}. The registry labels the problem \texttt{open}; this is a supplied status, not an independently verified mathematical assessment. Upstream revision: \texttt{d371ee917ac9647fbb06cd6fc529deec6f3dda82}.
\item \emph{The Kourovka Notebook (Unsolved Problems in Group Theory)}, original Problem 14.93, 14th issue (1999), proposed by I. B. Fesenko. The supplied background identifies the 21st issue edition (2026), \texttt{arXiv:1401.0300v48}; the collection description names E. I. Khukhro and V. D. Mazurov as editors.
\item Archive Problem 12.24 of the notebook is the supplied cross-reference for the definition of the Nottingham group, also called the ``Wild group''. Its text was not supplied.
\item M. V. Ershov, \emph{J. London Math. Soc.} \textbf{71} (2005), 362--378. The source's comment dated 2009 cites this reference as progress for \(p>2\). No article title or account of its results was supplied. These bibliographic and progress statements have not been independently verified.
\end{itemize}
\end{document}
